\documentclass[12pt]{article}
\usepackage{amsfonts}
\usepackage{geometry}
\usepackage{amsmath,amsthm,amssymb}
\usepackage{graphicx}
\usepackage[utf8]{inputenc}
\usepackage{setspace}
\usepackage{xcolor}
\usepackage{hyperref}
\usepackage{parskip}
\usepackage[hang, flushmargin, bottom]{footmisc}

\geometry{verbose,tmargin=1in,bmargin=1in,lmargin=1in,rmargin=1in,nomarginpar}
\hypersetup{breaklinks=true,hypertexnames=false,colorlinks=true,citecolor=teal,urlcolor=teal}
\setlength{\parskip}{0.8em}

\begin{document}

\begin{center}
    {\large \textbf{Application: Visiting Scholars Program}}\\[0.4em]
    {\large Banco Central de Chile}\\[0.8em]
    {\normalsize Lucas Schmitz\footnote{Yale University. \texttt{lucas.schmitz@yale.edu}}}
\end{center}

\bigskip

\noindent\textbf{Research Project}

\medskip

\noindent\textit{Competition and Credit Information in Consumer Lending Markets: Effects of Chile's REDEC Reform (Law 21.680)}

\medskip

This project uses structural Industrial Organization methods to study how information asymmetries and switching costs shape competition, credit pricing, and borrower welfare in consumer lending markets, with a focus on the effects of Law 21.680, which established Chile's Consolidated Debt Registry (REDEC) in July 2024. Prior to this reform, only banks and the largest savings cooperatives (CACs) were required to report to---and permitted to access---the CMF's N\'{o}mina de Deudores, while non-bank credit providers (OCNBs)---including retail lenders, cajas de compensaci\'{o}n, smaller cooperatives, and factoring and leasing firms, which together account for over 60\% of consumer credit placements---were excluded from the registry. The Banco Central's Informes de Estabilidad Financiera for the second semester of 2018 and the first semester of 2019 documented the risks arising from this fragmented information system, showing that the problem REDEC addresses had been a concern of the Bank well before the law's enactment. REDEC extends both the reporting obligation and access to the registry to all credit providers, changing the information available to lenders when evaluating borrowers. This creates a clean before-and-after comparison to estimate how broader information sharing affects credit pricing, competition between bank and non-bank lenders, and household over-indebtedness. Using the administrative credit records held at the Banco Central, I would estimate these effects and their distributional implications across borrower segments.

\bigskip

\noindent\textbf{Preferred Dates and Duration}

\medskip

\noindent I am available for a visit between June 20 and July 31, 2026, and I am happy to agree on a specific start date and duration within that window based on the Bank's needs and schedule.

\end{document}